\documentclass[10pt]{article}
\usepackage{amsmath,amsfonts,amssymb,mathtools,mathrsfs,bm,bbm,array,xcolor,geometry}
\allowdisplaybreaks
\setlength{\parindent}{0pt}
\setlength{\parskip}{4pt}
\emergencystretch=4em

\geometry{paperwidth=215mm,paperheight=290mm,hmargin=15mm,vmargin=16mm}
\begin{document}
\section*{P0177 Truck and Delivery Robot Rendezvous Planning}
\subsection*{Mathematical formulation}
\[
 T_o=\sum_{(a,b)\in P^T_o}d^T_{ab}
       +t_s(|H_o|+\mathbf1_{\{e_o\ne v_0\}}),\qquad
 R_o=\begin{cases}
 \sum_{(a,b)\in P^R_o}d^R_{ab}+|J_o|t_s+t_l+t_u,&J_o\ne\varnothing,\\
 0,&J_o=\varnothing,
 \end{cases}\quad c_o=\max(T_o,R_o).
\]

\begin{align*}
\min\quad &\sum_{o\in O}c_ox_o\\
\text{s.t.}\quad
&\sum_{o\in O(v)}x_o\ge1 &&v\in V,\\
&\sum_{o\in O^-(v)}x_o\le y_v &&v\in V,\\
&\sum_{o\in O^-(v)}x_o=\sum_{o\in O^+(v)}x_o &&v\in V,\\
&\sum_{o\in O^+(S)}x_o\ge y_v
 &&\varnothing\ne S\subseteq V\setminus\{v_0\},\ v\in S,\\
&\sum_{o\in O^+(v_0)}x_o\ge1,\qquad y_{v_0}=1,\\
&x_o,y_v\in\{0,1\}.
\end{align*}

\subsection*{Definitions and model scope}
A truck and one delivery robot leave a depot $v_0$ together, deliver parcels to customers $C$, and return together. Each customer has one parcel. The node set $V$ contains the depot and customers. Complete distances $D_{ab}$ are Euclidean distances from the given coordinates in kilometers. Truck and robot travel times in seconds are distance divided by their respective speeds $v_T,v_R$ in kilometers per hour, then multiplied by 3600. A customer requires service time $t_s$, while the depot has no customer service. Robot loading and unloading take times $t_l,t_u$, and the robot can carry at most parcel capacity $p$.

A joint operation $o$ has a common departure node $s_o$ and reunion node $e_o$, an ordered internal truck customer sequence $H_o$, and an ordered robot customer sequence $J_o$. Internal customers do not repeat or overlap between these two sequences, and the robot sequence does not exceed parcel capacity. All legal such operations belong to the available operation collection $O$. With no robot customers, the robot rides on the truck. Truck duration sums travel along consecutive links from departure through internal truck customers to reunion, plus customer service for each internal customer and for the reunion when it is a customer. Robot duration is zero for an empty robot sequence; otherwise it sums travel through its sequence to reunion, service for its own customers, and loading and unloading. The operation duration $c_o$ is the larger of truck and robot durations. Reunion customer service belongs to the arriving operation.

Zero-or-one operation choices $x_o$ and zero-or-one common-node records $y_v$ describe the plan. Every node, including the depot, must be visited by at least one selected operation; this is a coverage lower bound and does not impose exactly one visit across operations. For every node, the number of selected operations starting there cannot exceed its common-node record, and equals the number ending there. At least one selected operation ends at the depot and its common-node record is one. Therefore selected boundary operations balance at each node, with at most one starting there.

For every nonempty collection of nodes excluding the depot, and for each node within that collection, the number of selected operations entering the collection from outside must be at least that node's common-node record. This prevents an active group of boundary nodes from being isolated from the depot. The operator seeks the smallest sum of selected operation durations, each counted once. No arbitrary limit on the length of a forward customer sequence or on the set of available operations is imposed.

\subsection*{Inline notation and relations}
$V=\{v_0\}\cup C$; $v_0$; $d^T_{ab}=3600D_{ab}/v_T$; $d^R_{ab}=3600D_{ab}/v_R$; $t_s$; $t_l,t_u$; $p$; $o=(s_o,H_o,J_o,e_o)$; $H_o$; $J_o$; $|J_o|\le p$; $O$; $P^T_o=(s_o,H_o,e_o)$; $P^R_o=(s_o,J_o,e_o)$; $x_o$; $y_v$; $O(v)$; $v$; $O^-(v)$; $O^+(v)$; $O^+(S)$; $V\setminus S$; $S$; $\le1$; $\sum_{v\in S}y_v$.
\end{document}